% TOP_PROBLEM: {"id": 2928, "title": "Kirby Problem 4.52", "category_id": 7, "category": {"id": 7, "name": "topology", "display_name": "Topology", "description": "Properties preserved under continuous deformations.", "slug": "topology", "order_index": 7, "created_at": "2026-07-31T15:26:25.671Z"}, "status": "open"}
% FIRST_POSTED: null
% ENTRY_KIND: statement_only
% Imported from: batch13.tex; UnsolvedMath ID 2928
\documentclass[11pt]{article}
\usepackage[margin=25mm]{geometry}
\usepackage{fontspec}
\setmainfont{Latin Modern Roman}
\usepackage{amsmath,amssymb,mathtools}
\usepackage{unicode-math}
\setmathfont{Latin Modern Math}
\usepackage{xurl,hyperref}
\hypersetup{colorlinks=true,urlcolor=blue}
\setcounter{secnumdepth}{0}
\setlength{\parindent}{0pt}
\setlength{\parskip}{5pt}
\setlength{\emergencystretch}{3em}
\newcommand{\sref}[2]{\hyperlink{source:#1:#2}{[S#2]}}
\newcommand{\cataloguescope}[1]{\paragraph{Recorded target.}#1}
\begin{document}
% UnsolvedMath ID 2928; source catalogue code KP-4.52
\setcounter{section}{2927}
\section{Kirby Problem 4.52}\label{2928}
{\small\sffamily UnsolvedMath ID 2928 · Topology}
\subsection{Definitions and mathematical statement}

\paragraph{Shared notation.}$\mathbb N=\{1,2,\ldots\}$, and $\mathbb Z,\mathbb Q,\mathbb R,\mathbb C$ denote the integers, rationals, reals and complex numbers. Any other conventions are specified locally.


Let $M,N$ be closed connected orientable topological four-manifolds with isomorphic fundamental groups, identified with a good group $\pi$.

A group $\pi$ is good in the Freedman--Quinn disk-embedding sense. The accepted capped-grope criterion is: for every based height-$1.5$ disk-like capped grope $G$ and homomorphism $\rho:\pi_1(G)\to\pi$, there is an immersed disk in $G$ with the same framed attaching boundary whose double-point loops all map to $1$ under $\rho$. A capped grope is the framed four-dimensional thickening of a surface complex built by attaching higher-stage surfaces along symplectic-basis curves and capping its tips by disks; height $1.5$ means a first stage and a second stage on one half of a symplectic basis. A double-point loop follows the two preimage branches of a self-intersection, with paths to the basepoint. The source gives this criterion in Problem 4.46, Remarks (1)--(2).

A simple homotopy equivalence is a homotopy equivalence with zero Whitehead torsion. Stable homeomorphism means that, for some nonnegative integers $r,s$, $M\#r(S^2\times S^2)$ and $N\#s(S^2\times S^2)$ are homeomorphic. Choose compatible orientations. The simple topological structure set $\mathcal S^s_{\mathrm{TOP}}(M)$ consists of simple homotopy equivalences $f:P\to M$ preserving these orientations, modulo homeomorphisms of $P$ commuting with the markings up to homotopy. The group $L_5^s(\mathbb Z[\pi])$, with involution $g\mapsto g^{-1}$, acts by topological Wall realization of five-dimensional normal surgery data. Its identity structure is $[M,\mathrm{Id}_M]$.
\paragraph{Statement recorded in the supplied source.}
If $M$ and $N$ are simple homotopy equivalent and stably homeomorphic, must there exist a simple homotopy equivalence $f:N\to M$ and an element $A\in L_5^s(\mathbb Z[\pi])$ such that
\[
[N,f]=A\cdot[M,\mathrm{Id}_M]\quad\text{in }\mathcal S^s_{\mathrm{TOP}}(M)?
\]
The question seeks some marking on the fixed manifold $N$ in the identity orbit; it does not assert that every prechosen simple equivalence lies in that orbit. The source further asks whether there exist closed spin topological four-manifolds that are simple homotopy equivalent but not stably homeomorphic. Here spin means that the oriented stable topological tangent microbundle has vanishing second Stiefel--Whitney class $w_2$; the subsidiary existence question does not impose an additional good-group restriction. \sref{2928}{2}
\subsection{Short English statement}
For good fundamental groups, does stable homeomorphism plus simple homotopy equivalence imply a suitable Wall realization marking? Also, can closed spin four-manifolds be simply homotopy equivalent without being stably homeomorphic?
\subsection{Sources}
\begin{description}
\item[\hypertarget{source:2928:1}{S1}] ULAM.AI and the UnsolvedMath Contributors, UnsolvedMath: A Curated Collection of Open Mathematics Problems, record 2928 (KP-4.52). CC BY 4.0. \url{https://huggingface.co/datasets/ulamai/UnsolvedMath}.
\item[\hypertarget{source:2928:2}{S2}] R. İnanç Baykur, Robion C. Kirby and Daniel Ruberman, K3—A New Problem List in Low-Dimensional Topology, Mathematical Surveys and Monographs 295 (2026), author version, Problem 4.52 and Remarks, PDF page 231. \url{https://math.berkeley.edu/sites/default/files/surv-295-ruberman-watermarked-author-pdf.pdf}.
\end{description}
\label{end:2928}
\end{document}
